\documentclass[10pt,a4paper]{amsart}
\usepackage[margin=19mm]{geometry}
\usepackage{amsmath,amssymb,mathtools}
\usepackage[scr=rsfs,cal=euler]{mathalfa}
\usepackage{xfrac}
\usepackage[unicode,hidelinks,bookmarks=false]{hyperref}
\newcommand{\ie}{i.e.\ }
\newcommand{\cf}{cf.\ }
\newcommand{\resp}{respectively\ }
\newcommand{\Subsection}{\subsection}
\providecommand{\nobreakdash}{\nobreak\hbox{-}}
\setlength{\emergencystretch}{2em}
\multlinegap0pt
\usepackage{bm}
\usepackage{bbm}
\usepackage{dsfont}
\usepackage{xspace}
\usepackage{cancel}
\usepackage{mathtools}
\usepackage{amsthm}
\makeatletter
\@ifundefined{lem}{\newtheorem{lem}{Lemma}}{}
\makeatother
\newcommand{\ol}{\overline}
\title[Bruck conjecture for solutions of first-order partial differ]{Bruck conjecture for solutions of first-order partial differential equations in Cm}
\author{Sujoy Majumder, Nabadwip Sarkar, Debabrata Pramanik}
\date{Source: arXiv:2509.02576v2 (CC BY 4.0)}
\begin{document}
\maketitle

\noindent\emph{Source and licence.} Bruck conjecture for solutions of first-order partial differential equations in Cm. Version arXiv:2509.02576v2 (CC BY 4.0), distributed under the Creative Commons Attribution 4.0 International licence, as recorded in the arXiv licence metadata for this exact version and shown on the abstract page https://arxiv.org/abs/2509.02576v2. Full rights notice: licenses/arxiv-2509.02576-CC-BY-4.0.txt.\par\smallskip

\noindent\textit{Editorial scope.} Counted unit: the Lem whose source label is \texttt{t2.1} in the article (source0.tex lines 415--420, source label \texttt{t2.1}), delivered together with the article's own complete original proof of it (source0.tex lines 421--487, one \texttt{proof} environment of 67 lines). No mathematics is written, repaired or re-derived: statement and proof are the article's own text line for line. Required context retained uncounted: none. Every definition, display and label of those lines is retained unchanged. Omitted from this package and not redistributed: 1--414, 488--914. Nothing inside a retained excerpt is omitted or altered: every retained body is delivered exactly as the article's lines are written, so no omission receipt of any kind accompanies this package. Citations: the retained excerpts carry no citation command at all, so no bibliography entry of the article is transcribed and no reference is redistributed. Reference adaptation: none was needed and none was made; every reference inside the delivered unit resolves inside it, so no unresolved reference or citation is produced. Printed slips kept literal: no printed word, symbol, hypothesis or number of the counted statement or of its proof was altered. Layout and class: the article's own class is \texttt{amsart} with packages including indentfirst, amssymb, amsmath, amsthm, mathrsfs, setspace, enumerate, hyperref, cite; the wrapper is the lane's pinned \texttt{amsart} editorial wrapper at 10pt on A4, which restates the theorem-like environments the retained text needs and adds the article's own macro definitions that the retained text needs (\texttt{\textbackslash ol}), with extra wrapper packages (bm, bbm, dsfont, xspace, cancel, mathtools). The delivered numbering is the unit's own. Assets: no retained span contains a figure environment, a graphics-inclusion command, a PDF-page inclusion, a table or a TikZ picture; no image file is copied into this package, the package declares no asset dependency and no image file is redistributed with it. Licence: version arXiv:2509.02576v2 of this article is distributed under the Creative Commons Attribution 4.0 International licence, as recorded in the arXiv licence metadata for this exact version and shown on the abstract page https://arxiv.org/abs/2509.02576v2. A full rights notice is delivered at licenses/arxiv-2509.02576-CC-BY-4.0.txt. Characters: every retained excerpt is pure ASCII, so the delivered engine, XeLaTeX with its default Latin Modern text font, reports no missing character.
\begin{lem}\label{t2.1} Suppose that $f$ is a holomorphic function in $\ol{\mathbb{D}(0;R)}\;(0<R<+\infty)$. Then
\begin{align*}
 M(r,f)\leq \frac{2r}{R-r}A(R,f)+\frac{R+r}{R-r} |f(0)|
\end{align*}
holds for $0\leq \sqrt{m}r< R$.
\end{lem}

\begin{proof} We consider the following three cases.
\vspace{1.2mm}

{\bf Case 1.} Assume that $f$ is a constant. Let $f(z)=\alpha+i\beta$, where $\alpha$ and $\beta$ are real constants. Clearly $|f(0)|=\sqrt{\alpha^{2}+\beta^{2}}$, $M(r,f)=\sqrt{\alpha^{2}+\beta^{2}}$ and $A(r,f)=\alpha$. Then
\begin{align*}
 \frac{2r}{R-r}A(R,f)+\frac{R+r}{R-r}|f(0)|-M(r,f)=\frac{2r}{R-r}\Big(\alpha+\sqrt{\alpha^{2}+\beta^{2}}\Big).
 \end{align*}

Since $\alpha+\sqrt{\alpha^{2}+\beta^{2}}\geq \alpha+|\alpha|\geq 0$, we get
\begin{align*}
M(r,f)\leq \frac{2r}{R-r}A(R,f)+\frac{R+r}{R-r}|f(0)|.
\end{align*}
\vspace{1.2mm}

{\bf Case 2.} Assume that $f$ is non-constant and $f(0)=0$. Clearly $A(0,f)=0=M(0,f)$. Since both $A(r,f)$ and $M(r,f)$ are steadily increasing functions of $r$ and so for $r>0$, we have $A(r,f)>0$ and $M(r,f)>0$. Let $f(z)=u(x_1,y_1,\ldots,x_m,y_m)+\iota v(x_1,y_1,\ldots,x_m,y_m)$. Clearly
\begin{align*}
 2A(R,f)-f(z)=(2A(R,f)-u)+\iota (-v)\;\;\text{and}\;\;\Re\{2A(R,f)-f(z)\}=2A(R,f)-u.
 \end{align*}

For $0<\sqrt{m}r<R$, we have $0<A(\sqrt{m}r,f)\leq A(R,f)$. Since $u\leq A(\sqrt{m}r,f)$, we have $u\leq A(R,f)$ and $u<2A(R,f)$. Consequently $A(R,f)-u\geq 0$ and $2A(R,f)-u>0$. Note that
\begin{align}\label{pd1} 
|2A(R,f)-f(z)|^{2}=(2A(R,f)-u)^{2}+v^{2}=4A(R,f)[A(R,f)-u]+u^{2}+v^{2}\geq u^{2}+v^{2}.
\end{align}

Let 
\begin{align}\label{pd2} 
\phi(z)=\frac{f(z)}{2A(R,f)-f(z)}.
\end{align}

Clearly $\phi(z)$ is holomorphic in $\ol{\mathbb{D}(0;R)}$ and $\phi(0)=0$. Therefore using (\ref{pd1}) to (\ref{pd2}), we get $|\phi(z)|\leq 1$, for all $z\in \ol{\mathbb{D}(0;R)}$. Then by Schwarz's Lemma, we have $|\phi(z)|\leq \frac{1.r}{R}$, i.e., 
\begin{align}\label{pd3}
|\phi(z)|\leq \frac{r}{R}
\end{align}
holds for all $z\in\mathbb{C}^m[r]$, where $\sqrt{m}r<R$. Now from (\ref{pd2}), we have 
\begin{align}\label{pd4} 
|f(z)|=\Big|\frac{2A(R,f)\phi(z)}{1+\phi(z)}\Big|\leq \frac{2A(R,f)|\phi(z)|}{1-|\phi(z)|}.
\end{align}

Therefore using (\ref{pd3}) to (\ref{pd4}), we have
\begin{align}\label{pd5} 
|f(z)|\leq \frac{2A(R,f)\frac{r}{R}}{1-\frac{r}{R}}=\frac{2r}{R-r}A(R,f)
\end{align}
for all $z\in\mathbb{C}^m[r]$, where $\sqrt{m}r<R$. Since $f(0)=0$, using maximum modulus theorem to (\ref{pd5}), we have 
\begin{align*}
 M(r,f)\leq \frac{2r}{R-r}A(R,f)+\frac{R+r}{R-r}|f(0)|
 \end{align*}
holds for $0\leq \sqrt{m}r<R$.\vspace{1.2mm}

{\bf Case 3.} Assume that $f$ is non-constant and $f(0)\not=0$. Let $\phi(z)=f(z)-f(0)$. Clearly $\phi(0)=0$. Using maximum modulus theorem to (\ref{pd5}) (replacing $f$ by $\phi$), we get
\begin{align}\label{pd6} 
\max\limits_{||z||=r}|\phi(z)|\leq \frac{2r}{R-r}\max\limits_{z\in \mathbb{D}\langle 0;R\rangle}\Re\{\phi(z)\},
\end{align}
where $\sqrt{m}r<R$. Now we see that
\begin{align*}
 \max\limits_{||z||=r}|\phi(z)|=\max\limits_{||z||=r}|f(z)-f(0)|\geq \max\limits_{||z||=r}|f(z)|-|f(0)|=M(r,f)-|f(0)|
 \end{align*}
and 
\begin{align*}
 \max\limits_{z\in \mathbb{D}\langle 0;R\rangle}\Re\{\phi(z)\}=\max\limits_{z\in \mathbb{D}\langle 0;R\rangle}\Re\{f(z)-f(0)\}\leq \max\limits_{z\in \mathbb{D}\langle 0;R\rangle}\Re\{f(z)\}+|f(0)|=A(R,f)+|f(0)|.
 \end{align*}

Then from (\ref{pd6}), we deduce that 
\begin{align*}
 M(r,f)\leq \frac{2r}{R-r}A(R,f)+\frac{R+r}{R-r}|f(0)|
 \end{align*}
holds for $0\leq \sqrt{m}r<R$.
\end{proof}

\end{document}
